\chapter{Теория}

\section{Ping (Ubuntu)}

\subsection{Описание}
Утилита будет отправлять запросы ICMP echo-Request указанному узлу и фиксировать
ответы ICMP echo-Reply. Ping не определяет скорость передачи данных, а проверяет наличие
связности. Благодаря отображению времени между отправкой запроса и получения ответа
(RTT – Round Trip Time), можно косвенно определить загруженность сети передачи данных. \\

\subsection{Ключи}

С помощью различных ключей для команды, можно вывести дополнительную информацию
или убрать ненужную. \\

Ключ -n позволяет выводить информацию без расшифровки имени адреса, только IP
адрес. \\

С помощью ключа -n можно задать время работы команды в секундах, по истечению
которого она прекратит отправку пакетов.

\subsection{Результат работы}
\begin{enumerate}
    \item Размер эхо-запроса стандартные 64 байта (20 байт составляет IP-заголовок).
    \item Адрес назначения.
    \item Порядковый номер запроса.
    \item Количество переходов, которые может сделать запрос.
    \item Среднее время ответа от заданного узла.
    \item Общий вывод информации, в который входит количество отправленных, принятых пакетов, результат потери пакетов в процентах, время работы процесса и информация о
RTT (круговая задержка).
\end{enumerate}

RTT (Round Trip Time) – это время, необходимое для отправки пакета данных до
получателя и обратно к отправителю. По умолчанию команда traceroute отправляет три пакета
до каждого узла и выводит нам время, затраченное на отправку и возврат каждого из пакетов.
Выводится именно три параметра времени, так как по нескольким значениям можно сделать
более достоверный вывод о времени задержки в сети.

\section{Анализ в Wireshark}

\section{Практическое применение}

\begin{enumerate}
    \item Проверить доступность устройств.
    \item Анализ времени ответа.
    \item Проверка потери пакетов.
\end{enumerate}

\section{Traceroute (Ubuntu)}

\subsection{Описание}
Утилита Traceroute позволяет проследить маршрут следования данных до адреса
назначения.

\subsection{Установка}

\begin{lstlisting}
     sudo apt install inetutils-traceroute -y
\end{lstlisting}

\subsection{Результат работы}

\begin{enumerate}
    \item Порядковый номер маршрутизатора (узла).
    \item IP-адрес узла (в некоторых случаях с именем узла).
    \item RTT (круговая задержка).
\end{enumerate}

\subsection{Анализ в Wireshark}

\subsection{Практическое применение}
Используя команду traceroute можно отследить маршрут передачи данных, и по параметру
круговой задержки (RTT) найти узел, который вносит наибольшие временные задержки в
передачу данных. Это может быть по причине высокой нагрузки на сеть или из-запроблем с
оборудованием.
Также данная утилита может помочь отследить неправильную маршрутизацию. Например,
был настроен статический маршрут передачи данных между двумя узлами, но доступ к
искомому узлу не осуществляется. С помощью команды traceroute можно отследить узел,
после которого маршрут теряется. Это может быть, к примеру, маршрутизатор, в котором не
настроен маршрут до искомой сети. То есть, статический маршрут между узлами настроили
правильно, но один из маршрутизаторов на этом пути «не знает» сеть, в которую ему нужно
передать данные. После добавления данной сети в таблицу маршрутизации устройства,
проблема устранится (в рамках данного простого примера).



\endinput